\section{Related Work}
\noindent\textbf{Vision-language-action robot foundation policies.}
Recent robot foundation policies have rapidly scaled vision-language-action models for general manipulation through large-scale robot data, vision-language pretraining, diffusion or flow-matching action decoders, and cross-embodiment training~\cite{intelligence2025pi05visionlanguageactionmodelopenworld,nvidia2025gr00tn1openfoundation,liu2025rdt1bdiffusionfoundationmodel,wang2026qwenvlaunifyingvisionlanguageactionmodeling}.
While these models improve semantic grounding and language-conditioned action generation, they remain primarily driven by visual observations and action-level supervision.
In contact-rich dexterous manipulation, key physical states such as slip, grasp stability, contact force, and insertion progress are difficult to infer from vision alone.
TouchWorld extends this line of work by introducing tactile-conditioned action generation and high-frequency tactile refinement into a hierarchical VLA framework.

\noindent\textbf{Tactile representation learning and tactile manipulation policies.}
Tactile sensing provides direct measurements of local contact states, including pressure, deformation, slip, and hand-object interaction, which are difficult to infer reliably from RGB observations alone.
High-resolution tactile sensors such as GelSight~\cite{yuan2017gelsight} and DIGIT~\cite{lambeta2020digit} have enabled dense contact perception, while prior work has studied visuo-tactile prediction~\cite{li2019visgel}, visual pressure estimation~\cite{yang2022pressurevision,grady2024pressurevision++}, and tactile estimation from egocentric human interaction~\cite{zhou2026touchanythingdatasetframeworkbimanual}.
Building on these tactile representations, recent tactile manipulation policies incorporate touch into robot action models for contact-rich manipulation.
FTP-1 scales tactile policy learning across heterogeneous tactile sensors and contact-rich tasks~\cite{yuan2026ftp1generalistfoundationtactile}, while T-Rex studies tactile-reactive dexterous manipulation by exploiting high-frequency tactile signals for reactive control~\cite{niu2026trextactilereactivedexterousmanipulation}.
These works demonstrate that tactile observations are valuable not only for perception but also for contact-aware action generation.
However, tactile observations are often fused into a monolithic policy as an additional modality, or used primarily as reactive cues within the action-generation loop.
TouchWorld instead assigns touch two complementary roles: a predictive signal for visual-tactile goal generation and a fast feedback signal for high-frequency residual refinement.
This separation allows the visuo-tactile policy to preserve semantic and geometric progress while the tactile refinement layer performs online contact adaptation.

\noindent\textbf{Hierarchical, predictive, and reactive manipulation policies.}
Long-horizon manipulation often benefits from separating high-level task structure from low-level control.
Prior work has explored action chunking and diffusion-based visuomotor policies to generate temporally coherent short-horizon action sequences rather than single-step commands~\cite{zhao2023learningfinegrainedbimanualmanipulation,chi2025diffusion,black2025realtimeexecutionactionchunking}.
Hierarchical robot policies further use intermediate representations such as subtasks, paths, or visual goals to bridge language-level reasoning and local execution~\cite{li2025hamsterhierarchicalactionmodels}.
Predictive policies and robot world models extend this idea by forecasting future visual states, goal trajectories, or latent dynamics to guide goal-conditioned action generation~\cite{zhou2025act2goalworldmodelgeneral,zang2026tacforesightforceguidedtactileworld}.
In parallel, reactive tactile policies and residual controllers use recent tactile or force feedback to correct local execution errors during contact-rich manipulation~\cite{xue2025reactivediffusionpolicyslowfast,zhao2025touchbeginsvisionends}.
However, these directions are often studied separately: hierarchical and predictive policies mainly rely on language or visual goals, while reactive tactile policies focus on local contact adaptation without explicit long-horizon semantic structure.
TouchWorld combines these ideas in a unified tactile VLA framework by coupling executable subtask planning, predictive visual-tactile goals, nominal action-chunk generation, and online tactile residual refinement.


\section{Implementation Details}

\noindent\textbf{High-Level Planning Layer.}
The High-Level Planning Layer contains the Subtask Planner and the Tactile World Model.
We instantiate the Subtask Planner from Qwen3-VL-4B-Instruct and adapt it with supervised LoRA fine-tuning.
The Subtask Planner takes the task instruction, the current camera observation, and a compact memory of recent subtasks and execution states, and predicts the executable subtask used by the downstream policy.
The High-Level Planning Layer is updated at a fixed slow semantic rate in our experiments.
To make the Subtask Planner robust to stale or repeated histories, the SFT set includes teacher labels, oversampled rare phases, and noisy-memory variants.
The resulting Subtask Planner dataset contains 128,866 records sampled at this semantic update rate.
We use LoRA rank 16, alpha 32, dropout 0.05, a learning rate of $10^{-4}$, cosine decay with 0.1 warmup ratio, bfloat16 training, and 20 epochs.
At inference time, only the selected subtask is passed to the manipulation policy; intermediate reasoning is used only for logging and analysis.

\noindent\textbf{Visuo-Tactile Goal-Conditioned Policy.}
The nominal manipulation policy is implemented as a visuo-tactile diffusion Transformer.
It receives multi-view RGB observations, the current tactile observation rendered as an image, the predicted goal grid from the Tactile World Model when available, proprioceptive state, and a language prompt that concatenates the original task and the current subtask.
In this VLA branch, tactile input is represented uniformly as image-form context rather than by separate tactile encoders for different sensor types.
The model encodes visual, tactile-image, and language inputs as context tokens and denoises action tokens with a Transformer action expert.
On the Wuji platform, the policy predicts a 120-dimensional action vector over a 32-step action horizon.
The action vector consists of two 48-dimensional arm-hand action representations, a 9-dimensional head action, and 15 reserved dimensions.
Training uses the teleoperation action targets and the normalization statistics computed from the zarr training set.
We train the nominal policy for 30,000 optimization steps with global batch size 32, bfloat16 precision, AdamW, gradient clipping at 1.0, and a cosine learning-rate schedule with 1,000 warmup steps, peak learning rate $2.5\times10^{-5}$, and final learning rate $2.5\times10^{-6}$.

\noindent\textbf{Tactile World Model.}
The Tactile World Model is fine-tuned from Wan2.2-TI2V-5B and predicts short-horizon visual-tactile subgoal observations conditioned on the current multimodal context and selected subtask.
Each training sample is formatted as a 2-by-2 observation grid containing external RGB views and tactile-image observations, together with the task and subtask text.
The Tactile World Model uses large-scale human interaction video data and 10 hours of robot demonstration data.
At 30 FPS, the robot demonstrations correspond to approximately 1.08 million robot frames.
Robot training samples use 17-frame clips at $384 \times 224$ resolution.
The world model is adapted with LoRA on the DiT attention and feed-forward projections, using target modules $\{q,k,v,o,\mathrm{ffn}.0,\mathrm{ffn}.2\}$, rank 64, learning rate $10^{-4}$, and 50 epochs.
During deployment, the High-Level Planning Layer is still updated at the slow semantic rate, but the Tactile World Model is refreshed only when the high-level memory indicates a meaningful subtask or task-state change.
If the current phase is unchanged, the system reuses the previous predicted goal, which avoids repeatedly invoking the world model during stable execution phases.

\noindent\textbf{Tactile-Conditioned Refinement Policy.}
The refinement layer is implemented as TRT, a compact Tactile Residual Transformer over recent tactile history, proprioceptive state, sliding nominal-action lookahead windows, and VLA context tokens.
Unlike the VLA branch, this layer uses structured tactile histories and modality-specific tactile tokenization for different tactile signal types.
The implementation uses a short causal tactile window that captures recent contact changes.
TRT takes a sliding nominal-action lookahead window with $W=16$ and predicts a residual window over the same horizon.
The nominal VLA still produces the full 120-dimensional action, but the residual layer is trained only on a 58-dimensional tactile-sensitive action subspace covering the two wrist pose blocks and selected hand joints.
The predicted residual is scattered back into the 120-dimensional action vector, while the remaining action dimensions are kept from the nominal VLA output.
The residual Transformer uses $d_{\mathrm{model}}=512$, eight layers, eight attention heads, a $W$-step residual query set, and a small residual regularization weight of $10^{-4}$.
In residual training, the trained nominal VLA is attached to the refinement layer and kept frozen; the residual actor and tactile feedback encoder are optimized with masked MSE between the corrected action window and the demonstration action window.
We use AdamW with learning rate $10^{-4}$, weight decay $10^{-4}$, and the same action normalization as the nominal policy.

\noindent\textbf{Deployment schedule.}
The full system uses a multi-rate execution schedule.
The High-Level Planning Layer runs at a slow semantic rate.
The Tactile World Model is called only when the Subtask Planner changes the subtask or contact-relevant phase.
The visuo-tactile policy generates nominal action chunks with horizon $H=32$.
Within each chunk, TRT refines sliding nominal-action lookahead windows with $W=16$ at stride $C=4$ (e.g., offsets $\{0,4,8,12\}$) and commits the first $C$ corrected actions before refreshing the residual prediction.
For example, TRT conditions on nominal actions $[0,15]$ and commits corrected actions $[0,3]$, then conditions on $[4,19]$ and commits $[4,7]$.
The realized feedback rate depends on the robot control frequency, sensing pipeline, and commit interval $C$.
If the Subtask Planner or Tactile World Model is unavailable, the system falls back to the original task prompt and disables predicted-goal conditioning, while the visuo-tactile policy and tactile refinement layer remain executable.
